\documentclass[10pt,a4paper]{amsart}
\usepackage{amsmath,amssymb,amsthm,graphicx}
% Retained minimal preamble for the Stacks Project chapter more-algebra.tex (More on Algebra).
% The chapter file pinned-original-source.tex carries only \input{preamble} on its first line, so no
% preamble text lives inside the chapter itself.  The Stacks Project uses one single shared preamble for
% every chapter; the snapshot stages it identically as sources/stacks-chow/preamble.tex and
% sources/stacks-etale-cohomology/preamble.tex.  The block below is copied mechanically, line for line,
% from sources/stacks-chow/preamble.tex lines 14-18 (the Xy-pic setup) and 154-217 (the native theorem
% declarations together with the mathematical macros).  Nothing was added, removed or redefined; the
% wrapper class and the amsmath/amssymb/amsthm/graphicx base packages are supplied by the assembler.
% Chapter: more-algebra.tex -- More on Algebra.
% For commutative diagrams we use Xy-pic
\usepackage[all]{xy}

% We use 2cell for 2-commutative diagrams.
\xyoption{2cell}

% Theorem environments.
%
\theoremstyle{plain}
\newtheorem{theorem}[subsection]{Theorem}
\newtheorem{proposition}[subsection]{Proposition}
\newtheorem{lemma}[subsection]{Lemma}

\theoremstyle{definition}
\newtheorem{definition}[subsection]{Definition}
\newtheorem{example}[subsection]{Example}
\newtheorem{exercise}[subsection]{Exercise}
\newtheorem{situation}[subsection]{Situation}

\theoremstyle{remark}
\newtheorem{remark}[subsection]{Remark}
\newtheorem{remarks}[subsection]{Remarks}

\numberwithin{equation}{subsection}

% Macros
%
\def\lim{\mathop{\mathrm{lim}}\nolimits}
\def\colim{\mathop{\mathrm{colim}}\nolimits}
\def\Spec{\mathop{\mathrm{Spec}}}
\def\Hom{\mathop{\mathrm{Hom}}\nolimits}
\def\Ext{\mathop{\mathrm{Ext}}\nolimits}
\def\SheafHom{\mathop{\mathcal{H}\!\mathit{om}}\nolimits}
\def\SheafExt{\mathop{\mathcal{E}\!\mathit{xt}}\nolimits}
\def\Sch{\mathit{Sch}}
\def\Mor{\mathop{\mathrm{Mor}}\nolimits}
\def\Ob{\mathop{\mathrm{Ob}}\nolimits}
\def\Sh{\mathop{\mathit{Sh}}\nolimits}
\def\NL{\mathop{N\!L}\nolimits}
\def\CH{\mathop{\mathrm{CH}}\nolimits}
\def\proetale{{pro\text{-}\acute{e}tale}}
\def\etale{{\acute{e}tale}}
\def\QCoh{\mathit{QCoh}}
\def\Ker{\mathop{\mathrm{Ker}}}
\def\Im{\mathop{\mathrm{Im}}}
\def\Coker{\mathop{\mathrm{Coker}}}
\def\Coim{\mathop{\mathrm{Coim}}}

% Boxtimes
%
\DeclareMathSymbol{\boxtimes}{\mathbin}{AMSa}{"02}

%
% Macros for moduli stacks/spaces
%
\def\QCohstack{\mathcal{QC}\!\mathit{oh}}
\def\Cohstack{\mathcal{C}\!\mathit{oh}}
\def\Spacesstack{\mathcal{S}\!\mathit{paces}}
\def\Quotfunctor{\mathrm{Quot}}
\def\Hilbfunctor{\mathrm{Hilb}}
\def\Curvesstack{\mathcal{C}\!\mathit{urves}}
\def\Polarizedstack{\mathcal{P}\!\mathit{olarized}}
\def\Complexesstack{\mathcal{C}\!\mathit{omplexes}}
% \Pic is the operator that assigns to X its picard group, usage \Pic(X)
% \Picardstack_{X/B} denotes the Picard stack of X over B
% \Picardfunctor_{X/B} denotes the Picard functor of X over B
\def\Pic{\mathop{\mathrm{Pic}}\nolimits}
\def\Picardstack{\mathcal{P}\!\mathit{ic}}
\def\Picardfunctor{\mathrm{Pic}}
\def\Deformationcategory{\mathcal{D}\!\mathit{ef}}

% --- appended 2026-10-12 (Chief): the project own definitions of the reference/slogan/history
% display environments, copied mechanically from the shared Stacks preamble lines 8-12
% (\usepackage{verbatim} plus the three \newenvironment lines, which the chapter preamble keeps
% verbatim from the project preamble). Nothing added, removed or redefined: these environments
% carry no mathematical content and the project itself typesets them as comments. Without these
% three lines a counted statement whose own text uses them fails with Environment slogan undefined.
\usepackage{verbatim}
\newenvironment{reference}{\comment}{\endcomment}
\newenvironment{slogan}{\comment}{\endcomment}
\newenvironment{history}{\comment}{\endcomment}
\usepackage{cleveref}
\crefname{originallemma}{original lemma}{original lemmas}
\Crefname{originallemma}{Original lemma}{Original lemmas}
\makeatletter
\newlabel{algebra-lemma-separable-extension-preserves-reducedness}{{030U}{source}{}{sourcepointer-1}{}}
\newlabel{algebra-lemma-separable-extension-preserves-reducedness@cref}{{[originallemma][0][]030U}{[1][0][]source}{}{}{}}
\makeatother
\crefname{originallemma}{original lemma}{original lemmas}
\Crefname{originallemma}{Original lemma}{Original lemmas}
\makeatletter
\newlabel{fields-lemma-finite-purely-inseparable}{{09HI}{source}{}{sourcepointer-2}{}}
\newlabel{fields-lemma-finite-purely-inseparable@cref}{{[originallemma][0][]09HI}{[1][0][]source}{}{}{}}
\makeatother
\crefname{originallemma}{original lemma}{original lemmas}
\Crefname{originallemma}{Original lemma}{Original lemmas}
\makeatletter
\newlabel{lemma-construct-extension}{{35691}{source}{}{sourcepointer-3}{}}
\newlabel{lemma-construct-extension@cref}{{[originallemma][0][]35691}{[1][0][]source}{}{}{}}
\makeatother
\crefname{originallemma}{original lemma}{original lemmas}
\Crefname{originallemma}{Original lemma}{Original lemmas}
\makeatletter
\newlabel{lemma-extension-dvrs-formally-smooth}{{32989}{source}{}{sourcepointer-4}{}}
\newlabel{lemma-extension-dvrs-formally-smooth@cref}{{[originallemma][0][]32989}{[1][0][]source}{}{}{}}
\makeatother
\crefname{originallemma}{original lemma}{original lemmas}
\Crefname{originallemma}{Original lemma}{Original lemmas}
\makeatletter
\newlabel{lemma-nagata-goes-down}{{35667}{source}{}{sourcepointer-5}{}}
\newlabel{lemma-nagata-goes-down@cref}{{[originallemma][0][]35667}{[1][0][]source}{}{}{}}
\makeatother
\crefname{originalremark}{original remark}{original remarks}
\Crefname{originalremark}{Original remark}{Original remarks}
\makeatletter
\newlabel{remark-construction}{{33882}{source}{}{sourcepointer-6}{}}
\newlabel{remark-construction@cref}{{[originalremark][0][]33882}{[1][0][]source}{}{}{}}
\makeatother
\crefname{originalremark}{original remark}{original remarks}
\Crefname{originalremark}{Original remark}{Original remarks}
\makeatletter
\newlabel{remark-functoriality-construct-extension}{{35728}{source}{}{sourcepointer-7}{}}
\newlabel{remark-functoriality-construct-extension@cref}{{[originalremark][0][]35728}{[1][0][]source}{}{}{}}
\makeatother
\title{A solution containing a purely inseparable degree p step can be replaced by a separable solution}
\author{The Stacks Project Authors}
\date{stacks-project:more-algebra.tex:local-snapshot-20261012; GNU FDL 1.2 or later}
\begin{document}
\maketitle
\setcounter{subsection}{1}
\noindent\textit{Editorial note.} All delivered mathematical text is the author's own, copied from the pinned Stacks Project chapter file more-algebra.tex; no statement or proof step was rewritten, reordered or removed.\par\medskip
\noindent\textit{Editorial note.} The chapter file pinned-original-source.tex carries only \textbackslash{}input\{preamble\} on its first line, so no preamble text lives inside the chapter. The retained file sources/stacks-more-algebra/retained-preamble.tex is a mechanical, line-for-line copy of the project's single shared preamble as staged in the snapshot; nothing was added, removed or redefined and no package was added.\par\medskip
\noindent\textit{Editorial note.} The source URL is the source-level Stacks Project URL, as required when more than one card is produced from one source; the per-claim tag is carried by claim\_key.\par\medskip
\noindent\textit{Editorial note.} The author writes 'Some details omitted.' as the last sentence of the counted proof span; that sentence is kept verbatim and marks the step left to the reader, namely the verification of the two formal smoothness conditions on the pi\textasciicircum{}2-quotients. The references inside the counted proof span are recorded as non-load-bearing navigation pointers.\par\medskip
\noindent\textit{Allowed standard prerequisites (editorial).}\begin{itemize}
\item Solutions for an extension of discrete valuation rings and the construction of the rings A\_1, B\_1, L\_1. assumption (4) at statement line 35783 and used throughout the counted proof span. Supplies the hypothesis that is deformed into a separable solution.
\item Purely inseparable extensions of degree p in characteristic p and their description by x\textasciicircum{}p - a. assumption (5) at statement line 35784 and used at proof lines 35809-35813. Gives the element alpha and the constant a\_1 that the deformation replaces.
\item Nagata rings and finiteness of integral closures. assumption (3) at statement line 35782 and used at proof lines 35795-35803. Makes the auxiliary integral closures finite.
\item The extension-construction lemma and its functoriality, applied at f = pi\textasciicircum{}2. invoked at proof lines 35820-35822 and 35849-35852. Produces the integer n and the extensions A\_3 and B\_3 with the required isomorphism of quotients.
\item The criterion for formal smoothness of an extension of discrete valuation rings. invoked at proof lines 35880-35884. Is the final criterion that certifies K\_3/K\_1 as a solution.
\end{itemize}\medskip
\ifdefined\hypertarget\hypertarget{sourcepointer-1}{}\fi
\noindent\textit{Source reference.} Original lemma 030U: focused-seeds/AGP/sources/stacks-snapshot-20261011/tags.txt, source lines 3961--3961. Cross-chapter navigation pointer into the Stacks Project; the standard fact it supplies is applied inside the counted proof span and the argument resting on it stays there.\par\medskip
\ifdefined\hypertarget\hypertarget{sourcepointer-2}{}\fi
\noindent\textit{Source reference.} Original lemma 09HI: focused-seeds/AGP/sources/stacks-snapshot-20261011/tags.txt, source lines 12006--12006. Cross-chapter navigation pointer into the Stacks Project; the standard fact it supplies is applied inside the counted proof span and the argument resting on it stays there.\par\medskip
\ifdefined\hypertarget\hypertarget{sourcepointer-3}{}\fi
\noindent\textit{Source reference.} Original lemma 35691: sources/stacks-more-algebra/pinned-original-source.tex, source lines 35691--35691. Same-chapter navigation pointer; the statement it names is standard for the section and the argument that consumes it stays inside the counted proof span.\par\medskip
\ifdefined\hypertarget\hypertarget{sourcepointer-4}{}\fi
\noindent\textit{Source reference.} Original lemma 32989: sources/stacks-more-algebra/pinned-original-source.tex, source lines 32989--32989. Same-chapter navigation pointer; the statement it names is standard for the section and the argument that consumes it stays inside the counted proof span.\par\medskip
\ifdefined\hypertarget\hypertarget{sourcepointer-5}{}\fi
\noindent\textit{Source reference.} Original lemma 35667: sources/stacks-more-algebra/pinned-original-source.tex, source lines 35667--35667. Same-chapter navigation pointer; the statement it names is standard for the section and the argument that consumes it stays inside the counted proof span.\par\medskip
\ifdefined\hypertarget\hypertarget{sourcepointer-6}{}\fi
\noindent\textit{Source reference.} Original remark 33882: sources/stacks-more-algebra/pinned-original-source.tex, source lines 33882--33882. Same-chapter navigation pointer; the statement it names is standard for the section and the argument that consumes it stays inside the counted proof span.\par\medskip
\ifdefined\hypertarget\hypertarget{sourcepointer-7}{}\fi
\noindent\textit{Source reference.} Original remark 35728: sources/stacks-more-algebra/pinned-original-source.tex, source lines 35728--35728. Same-chapter navigation pointer; the statement it names is standard for the section and the argument that consumes it stays inside the counted proof span.\par\medskip
\begin{lemma}
\label{lemma-approximate-solution}
Let $p$ be a prime number. Let $A \subset B$ be an extension of discrete
valuation rings with fraction field extension $L/K$. Let $K_2/K_1/K$ be
a tower of finite field extensions. Assume
\begin{enumerate}
\item $K$ has characteristic $p$,
\item $L/K$ is separable,
\item $B$ is Nagata,
\item $K_2$ is a solution for $A \subset B$,
\item $K_2/K_1$ is purely inseparable of degree $p$.
\end{enumerate}
Then there exists a separable extension $K_3/K_1$
which is a solution for $A \subset B$.
\end{lemma}

\begin{proof}
Let us use notation as in Remark \ref{remark-construction}; we will use all
the observations made there. Since $L/K$ is separable, the algebra
$L_1 = L \otimes_K K_1$ is reduced
(Algebra, Lemma \ref{algebra-lemma-separable-extension-preserves-reducedness}).
Since $B$ is Nagata, the ring extension $B \subset B_1$ is finite
where $B_1$ is the integral closure of $B$ in $L_1$
and $B_1$ is a Nagata ring. Similarly, the ring $A$ is Nagata by
Lemma \ref{lemma-nagata-goes-down} hence $A \subset A_1$ is finite
and $A_1$ is a Nagata ring too. Moreover, the same assertions are true
for $K_2$, i.e., $L_2 = L \otimes_K K_2$ is reduced,
the ring extensions $A_1 \subset A_2$ and $B_1 \subset B_2$
are finite where $A_2$, resp.\ $B_2$ is the integral closure of
$A$, resp.\ $B$ in $K_2$, resp.\ $L_2$.

\medskip\noindent
Let $\pi \in A$ be a uniformizer. Observe that $\pi$ is a nonzerodivisor
on $K_1$, $K_2$, $A_1$, $A_2$, $L_1$, $L_2$, $B_1$, and $B_2$ and we have
$K_1 = (A_1)_\pi$, $K_2 = (A_2)_\pi$, $L_1 = (B_1)_\pi$, and $L_2 = (B_2)_\pi$.
We may write $K_2 = K_1(\alpha)$ where $\alpha^p = a_1 \in K_1$, see
Fields, Lemma \ref{fields-lemma-finite-purely-inseparable}.
After multiplying $\alpha$ by a power of $\pi$ we may and do assume
$a_1 \in A_1$. For the rest of the proof it is convenient to write
$K_2 = K_1[x]/(x^p - a_1)$ and $L_2 = L_1[x]/(x^p - a_1)$.
Consider the extensions of rings
$$
A'_2 = A_1[x]/(x^p - a_1) \subset A_2
\quad\text{and}\quad
B'_2 = B_1[x]/(x^p - a_1) \subset B_2
$$
We may apply Lemma \ref{lemma-construct-extension} to
$A'_2 \subset A_2$ and $f = \pi^2$ and to $B'_2 \subset B_2$ and $f = \pi^2$.
Choose an integer $n$ large enough which works for both of these.

\medskip\noindent
Consider the algebras
$$
K_3 = K_1[x]/(x^p - \pi^{2n} x - a_1)
\quad\text{and}\quad
L_3 = L_1[x]/(x^p - \pi^{2n} x - a_1)
$$
Observe that $K_3/K_1$ and $L_3/L_1$ are finite \'etale algebra extensions of
degree $p$. Consider the subrings
$$
A'_3 = A_1[x]/(x^p - \pi^n x - a_1)
\quad\text{and}\quad
B'_3 = B_1[x]/(x^p - \pi^n x - a_1)
$$
of $K_3 = (A'_2)_\pi$ and $L_3 = (B'_3)_\pi$. We are going to construct
a commutative diagram
$$
\xymatrix{
B'_2/\pi^{2n} B'_2 \ar[r]_{\psi'} & B'_3/\pi^{2n} B'_3 \\
A'_2/\pi^{2n} A'_2 \ar[r]^{\varphi'} \ar[u] & A'_3/\pi^{2n} A'_3 \ar[u]
}
$$
Namely, $\varphi'$ is the unique $A_1$-algebra isomorphism
sending the class of $x$ to the class of $x$.
Similarly, $\psi'$ is the unique $B_1$-algebra isomorphism sending
the class of $x$ to the class of $x$. By our choice of $n$ we obtain, via
Lemma \ref{lemma-construct-extension} and
Remark \ref{remark-functoriality-construct-extension}
finite ring extensions $A'_3 \subset A_3$ and $B'_3 \subset B_3$
such that $A'_3 \to B'_3$ extends to a ring map $A_3 \to B_3$
and a commutative diagram
$$
\xymatrix{
B_2/\pi^2 B_2 \ar[r]_\psi & B_3/\pi^2B_3 \\
A_2/\pi^2 A_2 \ar[r]^\varphi \ar[u] & A_3/\pi^2A_3 \ar[u]
}
$$
with all the properties asserted in the references mentioned above
(in particular $\varphi$ and $\psi$ are isomorphisms).

\medskip\noindent
With all of this data in hand, we can finish the proof. Namely, we
first observe that $A_3$ and $B_3$ are finite products of
Dedekind domains with $\pi$ contained in all of the maximal ideals.
Namely, if $\mathfrak p \subset A_3$ is a maximal ideal, then
$\pi \in \mathfrak p$ as $A \to A_3$ is finite. Then $\mathfrak p/\pi^2 A_3$
corresponds via $\varphi$ to a maximal ideal in $A_2 / \pi^2A_2$
which is principal as $A_2$ is a finite product of Dedekind domains.
We conclude that $\mathfrak p/\pi^2 A_3$ is principal and hence by Nakayama 
we see that $\mathfrak p (A_3)_\mathfrak p$ is principal.
The same argument works for $B_3$. We conclude that
$A_3$ is the integral closure of $A$ in $K_3$
and that $B_3$ is the integral closure of $B$ in $L_3$.
Let $\mathfrak q \subset B_3$ be a maximal ideal lying
over $\mathfrak p \subset A_3$. To finish the proof we have
to show that $(A_3)_\mathfrak p \to (B_3)_\mathfrak q$ is
formally smooth in the $\mathfrak q$-adic topology. By the criterion of
Lemma \ref{lemma-extension-dvrs-formally-smooth} it suffices to show that
$\mathfrak p (B_3)_\mathfrak q = \mathfrak q (B_3)_\mathfrak q$
and that the field extension $\kappa(\mathfrak q)/\kappa(\mathfrak p)$
is separable. This is true because we may check both assertions
by looking at the ring map $A_3/\pi^2 A_3 \to B_3/\pi^2 B_3$
and this is isomorphic to the ring map
$A_2/\pi^2 A_2 \to B_2/\pi^2 B_2$ where the corresponding
statement holds by our assumption that $K_2$
is a solution for $A \subset B$. Some details omitted.
\end{proof}

\end{document}
